\clearpage
\turntitle{3}{The War Spreads}
  {The first Anti-Guerrilla Operation drives Tito west; uprisings bring the
  Partisans to brigade strength.}

\begin{lessonbox}
Game-Turn 3 activates two major systems: Axis Anti-Guerrilla Operations
\rulesref{8.4} and the Uprising Segment \rulesref{7.47, 7.5}. The
starting-Zone restriction has also expired.
\end{lessonbox}

\begin{center}
\begin{minipage}[t]{0.48\linewidth}
  \centering
  \includegraphics[width=\linewidth]{playbook-turn-3-map.pdf}
  \mapcaption{Serbia: AGO deployment and the reaction route west}
\end{minipage}
\hfill
\begin{minipage}[t]{0.48\linewidth}
  \centering
  \includegraphics[width=\linewidth]{playbook-turn-3-west-map.pdf}
  \mapcaption{Bosnia and Croatia: Tito's refuge and uprising checks}
\end{minipage}
\end{center}

\begin{multicols}{2}
\small
\phase{A. Reinforcement and Collaboration}

Place German 113 in Bosnia and the 4th and 5th Ustaša brigades in Croatia.
Transfer German 125 as scheduled \rulesref{13.91}.

Resolve collaboration:

\begin{tabularx}{\linewidth}{@{}Xccc@{}}
  \toprule
  Stack & Roll & Mod. & Result \\
  \midrule
  Serbia neutral (2) & \die{5} & 0 & pro-Yugoslav \\
  Serbia pro-Yugoslav (1) & \die{3} & +1 & Neutral \\
  Montenegro pro-Axis (4) & \die{5} & $-1$ & Neutral \\
  \bottomrule
\end{tabularx}

The Axis Player now has no Chetnik groups under control, but this calm will not
last.

\phase{A. Axis Anti-Guerrilla Operations Phase}

\textbf{Planning.} The Axis declares one operation and secretly records
\mbox{\emph{Serbia}~\mapbadge{1}}. German 704, 714, 717, and 342 are selected; each
could reach Serbia in a normal Movement Phase and the total is below seven
division equivalents \rulesref{8.41--8.42}.

\textbf{Reaction.} The Yugoslav Player rolls \die{3}. Tito and two Partisan
groups leave the Serbia Mountain for the Bosnia Partisan
Hide-away~\mapbadge{2}. Tito counts as one of the three reacting units. Four groups
remain in the Serbia Mountain.

\columnbreak

\textbf{Deployment and combat.} Put 714 and 717 in the Axis Mountain and 704
and 342 in the Axis Hide-away:

\begin{itemize}[leftmargin=1.2em,itemsep=2pt]
  \item \textbf{Mountain:} 24 SP attacks 4 SP at 6:1. The AGO adjustment moves
    this to 8:1. A \die{1} produces D1, doubled to D2. Two groups are lost and
    two retreat to the Bosnia Partisan Hide-away.
  \item \textbf{Hide-away:} 30 SP attacks the lone group that retreated from
    Užice. Its defensive strength is doubled to 2, yielding 15:1, capped at
    12:1. A \die{1} produces De and eliminates it.
\end{itemize}

The four German divisions are removed pending redeployment
\rulesref{8.47}. The Yugoslav Player has lost three groups, but the reaction
saved Tito and gave the survivors a common refuge in Bosnia.

\begin{commentarybox}
An AGO changes the ratio by two columns, doubles numbered results, and sends
retreating Yugoslav units to an adjacent Zone. Axis units are never forced to
retreat during the operation \rulesref{8.46}.
\end{commentarybox}

\subsection{Situation After the Operation}

\begin{tabularx}{\linewidth}{@{}>{\bfseries}l X@{}}
  Bosnia Hide-away & Tito and 4 Partisan groups \\
  Serbia & 5 groups at Novi Pazar; no Partisans in the Zone Display \\
  Slovenia Mountain & 5 Partisan groups \\
  Montenegro Hide-away & 1 Partisan group \\
  Partisan total & 15 groups \\
\end{tabularx}
\end{multicols}

\clearpage
\turntitle{3}{From Groups to Brigades}
  {Recruitment and two under-garrisoned Zones fill the Partisan counter pool.}

\begin{multicols}{2}
\small
\phase{B. Yugoslav Player-Turn}

Move Tito and all four Bosnia groups from the Hide-away to the Partisan
Mountain, then send one group to Tuzla. Three groups remain with Tito. Keep the
five-group stack at Novi Pazar and the five groups in the Slovenia Mountain.
Move the lone Montenegro group back to Cetinje.

\phase{C. Victory Point Stage}

Novi Pazar, Tuzla, and Cetinje yield 4 VP. The Slovenia Mountain yields 1 VP
for its 5 SP. The Bosnia Mountain has only 3 SP and yields none. Gain
\textbf{5 VP}; the running total becomes \textbf{14}.

\phase{D. Axis Player-Turn}

The Axis emphasizes garrison duty and initiates no Objective combat. Units
used in the Serbia AGO remain off map until the Redeployment Phase, but may
count toward Serbia's garrison because that was the target Zone
\rulesref{7.54}.

\phase{E. Recruitment and Tito Segments}

The Partisans begin the Terminal Stage with 15 groups:

\begin{tabularx}{\linewidth}{@{}Xr@{}}
  \toprule
  Source & Groups \\
  \midrule
  Novi Pazar: \die{6} $\times$ one-half, cap 5 & 3 \\
  Slovenia Mountain: $5 \div 4$ & 1 \\
  Tuzla: \die{6} $\times$ one-half, cap 1 & 1 \\
  Cetinje: \die{4} $\times$ one-half, cap 1 & 1 \\
  Tito in Bosnia & 2 \\
  \midrule
  Total before uprisings & 23 \\
  \bottomrule
\end{tabularx}

\columnbreak

\phase{E. Uprising Segment}

\textbf{Croatia~\mapbadge{3}.} Two Italian divisions, the Croat Guard regiment,
and five Ustaša brigades equal $4\tfrac{5}{6}$ divisions. Drop the fraction
when calculating the shortfall: Croatia is four divisions below its Garrison
Value of 8. The uprising roll is \die{1}; with a modifier of 0, the modified
result is 1. Four short divisions times 1 creates \textbf{4 groups}. An
alignment roll of \die{4} is even, so all are Partisan in this P-aligned Zone.

\textbf{Bosnia~\mapbadge{4}.} German 718 and 113 provide two divisions against
a Garrison Value of 6. A \die{2}, less Bosnia's modifier of 1, gives a modified
result of 1. Four more Partisan groups would be created after an even alignment
roll of \die{2}, but only three Partisan group counters remain. Place those
three in the Bosnia Partisan Hide-away \rulesref{7.53, 7.68}.

The Partisans now have all \textbf{30 group counters} on the map and have
achieved brigade strength \rulesref{7.62A}.

\phase{E. Guerrilla Status Phase}

The first brigade must be created where Tito is located. Five groups occupy
the Bosnia Mountain with Tito: three moved there and two arrived in the Tito
Segment. Replace three with the 31st Partisan Brigade. Two groups remain with
Tito \rulesref{7.63, 7.67}.

\begin{commentarybox}[title={The important trigger}]
The first Partisan brigade has appeared. Beginning next Game-Turn, German
movement expands into the Italian sphere, the Bulgarian 25th and 27th
Divisions become available, and the Axis may attempt to identify Tito
\rulesref{6.64, 9.41, 13.91D}.
\end{commentarybox}
\end{multicols}

\turnledger{+5}{14}{27 groups + 1 brigade}{2 / 0 / 5}{1 of 8}

\vspace{0.7em}
\begin{multicols}{2}
\small
\subsection{Checkpoint Before Game-Turn 4}

\begin{tabularx}{\linewidth}{@{}>{\bfseries}l X@{}}
  Bosnia Mountain & Tito, 31st Brigade, 2 groups \\
  Bosnia Hide-away & 3 groups \\
  Tuzla & 2 groups \\
  Croatia Hide-away & 4 groups \\
  Slovenia Mountain & 6 groups \\
  Novi Pazar & 8 groups \\
  Cetinje & 2 groups \\
\end{tabularx}

Redeploy the four German AGO divisions in Serbia. The first-brigade trigger is
now permanent even if Partisan strength later falls below 30 group
equivalents.

\columnbreak

\subsection{What the Turn Demonstrated}

\begin{itemize}[leftmargin=1.2em,itemsep=3pt]
  \item An AGO is planned before reaction, so the Axis commits units without
    knowing exactly which guerrillas will remain.
  \item Reaction can preserve Tito and concentrate scattered survivors in an
    adjacent Zone.
  \item AGO participants still count toward the target Zone's garrison while
    awaiting redeployment.
  \item Objective recruitment, Tito, and uprisings all resolve before the
    Guerrilla Status Phase.
  \item Counter availability is a hard limit: excess uprising reinforcements
    are lost.
\end{itemize}

\begin{commentarybox}[title={A real decision}]
The first brigade need not be the 31st, and the Yugoslav Player could leave
all five Bosnia Mountain groups unorganized. Creating it now is stronger and
activates several Axis responses; delaying it conceals Tito for longer.
\end{commentarybox}
\end{multicols}
